\chapter{General usage and description}
\label{ch:general}

\section{What pyabr3 does}

pyabr3\index{pyabr3} records auditory brainstem responses\index{ABR}
(ABRs). It presents clicks and tone pips through a National Instruments
digital-to-analog converter and Tucker-Davis Technologies (TDT) attenuators,
records the amplified scalp potential with a TDT RP2.1 processor, and shows the
responses while the recording runs. Every repetition of every stimulus is saved in
its own file.

The program has three parts:

\begin{itemize}
  \item \textbf{Acquisition}, in the main window (Chapters~\ref{ch:acquisition}
        and~\ref{ch:online}).
  \item \textbf{Calibration} of the microphone and speaker, in its own window
        (Chapter~\ref{ch:calibration}). This replaces the calibration part of the
        MATLAB program, ABR4\index{ABR4}.
  \item \textbf{A data viewer} for recordings made by pyabr3 or ABR4
        (Chapter~\ref{ch:viewing}).
\end{itemize}

The data viewer is for checking recordings. Threshold measurement and other analysis
are done with separate code, which this manual does not cover.

\section{Example workflow}
\index{workflow}

Figure~\ref{fig:workflow} shows a typical recording day. The numbered steps in
Section~\ref{sec:example-session} follow the same order.

\begin{figure}[htbp]
\centering
\begin{tikzpicture}[
    node distance=5mm and 9mm,
    wfbox/.style={draw=uncblue!80!black,fill=uncblue!12,rounded corners=2pt,thick,
                 minimum width=46mm,minimum height=10mm,align=center,font=\small},
    decide/.style={wfbox,fill=orange!15,draw=orange!70!black},
    flow/.style={-{Latex[length=2.5mm]},thick,draw=black!70},
    lbl/.style={font=\scriptsize\sffamily\bfseries,text=black!60}]
  % left column: preparation
  \node[wfbox] (start) {Start pyabr3\\\footnotesize\texttt{uv run python main.py}};
  \node[wfbox,below=of start] (mic) {Calibrate microphone\\\footnotesize 94~dB calibrator};
  \node[wfbox,below=of mic] (spk) {Calibrate speaker};
  \node[wfbox,below=of spk] (chk) {Check speaker at 75~dB};
  \node[decide,below=of chk] (ok) {Check is flat at 75~dB?};
  \node[wfbox,below=of ok] (read) {\gui{Read Cal File}\\\footnotesize select the new calibration};
  % right column: recording
  \node[wfbox,right=of start] (animal) {Prepare animal, place\\electrodes and speaker};
  \node[wfbox,below=of animal] (subj) {Enter \gui{Subject Data}};
  \node[wfbox,below=of subj] (proto) {Choose \gui{Protocol}};
  \node[wfbox,below=of proto] (test) {\gui{Test Acquisition}\\\footnotesize nothing is saved};
  \node[wfbox,below=of test] (run) {\gui{Start Acquisition}\\\footnotesize watch the displays};
  \node[wfbox,below=of run] (file) {Data saved in a new\\serial-numbered folder};
  \node[wfbox,below=of file] (view) {Review in the data viewer};
  % arrows
  \draw[flow] (start) -- (mic);
  \draw[flow] (mic) -- (spk);
  \draw[flow] (spk) -- (chk);
  \draw[flow] (chk) -- (ok);
  \draw[flow] (ok) -- node[right,font=\scriptsize]{yes} (read);
  \draw[flow] (ok.west) -- ++(-5mm,0) node[left,font=\scriptsize,align=right]{no: find\\the problem}
              |- (spk.west);
  \draw[flow] (read.east) -- ++(4.5mm,0) |- (animal.west);
  \draw[flow] (animal) -- (subj);
  \draw[flow] (subj) -- (proto);
  \draw[flow] (proto) -- (test);
  \draw[flow] (test) -- (run);
  \draw[flow] (run) -- (file);
  \draw[flow] (file) -- (view);
  \draw[flow] (file.east) -- ++(6mm,0) |- node[pos=0.25,right,font=\scriptsize,align=left]{next\\protocol} (proto.east);
  % headings
  \node[lbl,above=2mm of start] {SOUND SYSTEM};
  \node[lbl,above=2mm of animal] {RECORDING};
\end{tikzpicture}
\caption{A typical recording day. Calibration is not repeated for every animal; see
Section~\ref{sec:cal-overview}.}
\label{fig:workflow}
\end{figure}

\section{Where the program can run}
\index{Windows}\index{macOS}

\begin{center}
\begin{tabular}{@{}lll@{}}
\toprule
Task & Windows, with the rig hardware & macOS \\
\midrule
Acquisition & Yes & No \\
Calibration & Yes & Simulated hardware only \\
Viewing calibrations & Yes & Yes \\
Data viewer & Yes & Yes \\
\bottomrule
\end{tabular}
\end{center}

\section{Installing and starting}
\index{installation}\index{starting the program}

pyabr3 requires Python 3.13. Dependencies are managed with \file{uv}\index{uv}.
On the rig computer, the NI-DAQmx\index{NI-DAQmx} driver and the TDT drivers must
also be installed.

Run these commands from the top directory of the repository. The program finds its
configuration and protocols relative to the working directory, and will not start
from anywhere else.

\begin{lstlisting}
uv sync                       # create or update the environment
uv run python main.py         # start pyabr3
uv run python main.py test    # start with config/abrs_test.cfg
\end{lstlisting}

Without an argument, \file{main.py} uses \file{config/abrs.cfg} on Windows and
\file{config/abrs\_test.cfg} on other systems.

The calibration window and the data viewer can also be started by themselves:

\begin{lstlisting}
uv run python -m src.calibration.calibration_gui          # on the rig
uv run python -m src.calibration.calibration_gui --mock   # simulated hardware
uv run python -m src.viewer.data_viewer [file]
\end{lstlisting}

Stand-alone applications are built with \file{build\_windows.sh} and
\file{build\_macos.sh}\index{stand-alone application}.

\begin{caution}
The RP2.1 circuit is loaded from a fixed location,
\file{c:\textbackslash users\textbackslash experimenters\textbackslash desktop\textbackslash pyabr\textbackslash tdt\textbackslash abrs\_v2.rcx}.
The \key{rcx\_file} entry in the configuration file is not used. On a different
installation path, the program stops with a ``file not found'' message.
\end{caution}

\section{The main window}
\label{sec:main-window}
\index{main window}

The window has a column of parameters and controls on the left, and three tabs on the
right: \gui{Recording}, \gui{Preview} and \gui{Calibration}
(Figure~\ref{fig:main_window}). Chapter~\ref{ch:online} describes the tabs.

\figorplaceholder{main_window}{The main window before a recording, with the
\file{clickIO} protocol selected. This picture was taken on a computer with no
acquisition hardware.}{Main window at startup.}

\begin{center}
\begin{tabularx}{\textwidth}{@{}lX@{}}
\toprule
Group & Use \\
\midrule
\gui{Protocol} & Selects the protocol file. The list is read when the program starts. \\
\gui{Parameters} & Shows the values from the protocol file (see the caution below). \\
\gui{Subject Data} & Subject ID, age, sex, weight, strain, genotype and treatment group.
  These are saved with the data. \\
\gui{Actions} & The buttons that run the program (Table~\ref{tab:actions}). \\
\gui{Status} & Devices found, the current mode, and progress through the protocol. \\
\gui{Quit} & Stops the presentation and closes the program. \\
\bottomrule
\end{tabularx}
\end{center}

\begin{table}[htbp]
\centering
\caption{The \gui{Actions} group.}
\label{tab:actions}
\begin{tabularx}{\textwidth}{@{}lX@{}}
\toprule
Button & Action \\
\midrule
\gui{New Filename} & Has no effect on the names of saved files
  (Section~\ref{sec:data-files}). \\
\gui{Test Acquisition} & Presents the protocol and displays responses. Saves nothing. \\
\gui{Start Acquisition} & Presents the protocol, displays responses, and saves them. \\
\gui{Pause}, \gui{Resume} & Holds the presentation, and continues from the same place. \\
\gui{Stop} & Ends the presentation. It cannot be resumed. \\
\gui{Load Data File} & Opens the data viewer (Chapter~\ref{ch:viewing}). \\
\gui{Read Cal File} & Chooses the speaker calibration to use
  (Section~\ref{sec:cal-select}). \\
\gui{Calibrate} & Opens the calibration window (Chapter~\ref{ch:calibration}). \\
\bottomrule
\end{tabularx}
\end{table}

The \gui{Parameters} group\index{parameters!not editable} shows the values in the
selected protocol file and cannot be edited, so that a protocol of a given name is the
same from subject to subject. To use different values, save them as a new protocol file
(Section~\ref{sec:protocol-files}).

\gui{Start Acquisition} will not begin until a \gui{Subject ID}\index{Subject ID} has been
entered. The ID becomes part of the data folder name, so it cannot contain any of the
characters \verb+\ / : * ? " < > |+. \gui{Test Acquisition} saves nothing, and needs no
ID. Recordings made before this requirement store the word ``testing'' as their ID.

\section{Example session: a click intensity series}
\label{sec:example-session}
\index{example session}

\begin{enumerate}
  \item Turn on the hardware (Chapter~\ref{ch:wiring}) and start pyabr3. Check that
        \gui{Status}~/~\gui{Devices} lists the NI card, the RP2.1 and the PA5.
  \item If the sound system needs calibrating, follow
        Section~\ref{sec:cal-procedure}.
  \item Press \gui{Read Cal File} and select the calibration to use. Check its date
        and shape in the \gui{Calibration} tab.
  \item Prepare the animal, and place the electrodes and speaker.
  \item Fill in \gui{Subject Data}, including the \gui{Subject ID}, and choose the
        \gui{Data Path}.
  \item Select \file{click\_test.cfg} in \gui{Protocol}, and press
        \gui{Test Acquisition}. This presents 10 sweeps of clicks at 94~dB SPL. Check
        for a clean response in the \gui{Recording} tab, then press \gui{Stop}.
  \item Select \file{clickIO.cfg} and press \gui{Start Acquisition}. This protocol
        runs 17 levels from 90 to 10~dB SPL and takes about 7 minutes.
  \item The recording is complete when \gui{Status}~/~\gui{Mode} shows Finished.
  \item Open one of the files with \gui{Load Data File} to review the recording.
\end{enumerate}

\section{The configuration file}
\label{sec:config}
\index{configuration file}

The configuration files are in \file{config/}. The rig uses \file{abrs.cfg}.

\begin{center}
\begin{tabularx}{\textwidth}{@{}lXl@{}}
\toprule
Key & Meaning & Used? \\
\midrule
\key{required\_hardware} & Devices to connect: \texttt{NIDAQ}, \texttt{RP21},
  \texttt{PA5} & Yes \\
\key{NI\_device} & NI device name, from the NI-DAQmx utility & Yes \\
\key{calfile} & Speaker calibration read at startup. Its folder is also where the
  calibration window saves its files & Yes \\
\key{click\_reference\_dbspl} & SPL of a click at 0~dB attenuation (105.1) & Yes \\
\key{datapaths} & Named locations for saved recordings, one per line
  (Section~\ref{sec:data-files}). A single \key{datapath} entry also works & Yes \\
\key{rcx\_file} & RP2.1 circuit & No \\
\key{NI\_samplerate} & NI output rate & No (fixed at 500~kHz) \\
\key{RP21\_samplerate} & RP2.1 input rate & No (fixed at 24414.0625~Hz) \\
\key{soundcard\_output\_samplerate} & Sound card rate & No \\
\bottomrule
\end{tabularx}
\end{center}

\section{Protocol files}
\label{sec:protocol-files}
\index{protocol files}

Each file in \file{protocols/} with the extension \file{.cfg} is one protocol. To make
a new one, copy an existing file, edit it, and restart the program. This is
\file{clickIO.cfg}:

\begin{lstlisting}
# ABR Parameter file for click IO
# All times in seconds, freq in Hz
protocol:
    protocolname: 'clickIO'
    stimulustype: 'click'

stimuli:
    wave_duration: 1.2  # seconds
    stimulus_duration: 1e-4  # seconds
    stimulus_risefall: 5e-4  # seconds
    nstim: 40
    delay: 0.003 # seconds
    interval: 25e-3  # between clicks
    nreps: 20
    default_frequency: 4000.0
    default_spl: 80.0
    stimulus_period: 1e-2
    alternate: True
    dblist: [90.0, 85.0, 80.0, ... 15.0, 10.0]
    freqlist: []
recording:
    record_dur: 8e-3
    LPF: 3000.0
\end{lstlisting}

Chapter~\ref{ch:acquisition} explains each entry for each type of stimulus. The
\key{.toml} files in the same folder are not read by the program.

\section{Data files}
\label{sec:data-files}
\index{data files}\index{file names}

The \gui{Data Path}\index{Data Path} list holds the names in the \key{datapaths} entry
of the configuration file. A person can have several names, for example one for each
project. If there is more than one name, \gui{Start Acquisition} will not begin until
one has been chosen:

\begin{lstlisting}
datapaths:
    Reggie: "d:\\data\\reggie\\abr_data"
    Paul_projectA: "c:\\users\\experimenters\\desktop\\pyabr\\abr_data"
\end{lstlisting}

Each press of \gui{Start Acquisition} creates a new folder below the chosen location
(\file{<datapath>} below), and each repetition is written to one file in it:

\begin{lstlisting}
<datapath>/<date>-<Subject ID>_<serial>/<protocol>/<date>_<stimulus type>_<wave>_<repetition>.p
<datapath>/2024-12-10-p572_000/clickIO/2024-12-10_click_000_001.p
\end{lstlisting}

\file{<serial>} is the first number, from 000, for which that subject and day does not
already have a folder for the protocol. Different protocols on one subject therefore
share \file{\_000}, and a repeat of a protocol goes to \file{\_001}. Nothing is
overwritten.

\file{<wave>} counts from 0 and identifies the stimulus in the protocol.
\file{<repetition>} counts from 1. A file is a Python pickle\index{pickle} holding one
dictionary:

\begin{center}
\begin{tabularx}{\textwidth}{@{}lX@{}}
\toprule
Key & Contents \\
\midrule
\key{data} & The recorded trace, in volts at the amplifier output \\
\key{record\_frequency} & A/D sample rate, Hz \\
\key{subject\_data} & The \gui{Subject Data} fields \\
\key{protocol} & The protocol as used \\
\key{stimuli} & Levels, frequencies and attenuations of every stimulus. The waveforms
  themselves are not stored \\
\key{calibration} & The speaker calibration in use \\
\key{wave\_number}, \key{repetition} & Position in the protocol \\
\key{timestamp} & Time of writing \\
\bottomrule
\end{tabularx}
\end{center}

\gui{New Filename}\index{overwriting data} has no effect on where data are saved.

The folder arrangement below is the one that the analysis program \file{read\_abr.py}
expects. Recordings must be copied into it by hand before that program can use them:

\begin{lstlisting}
experiment_directory/
    CBA_M_N001_p48_NT/          # strain_sex_ID_age_...
        Clicks/   Tones/   Interleaved_plateau/
\end{lstlisting}

\needsinput{The lab's standard folder convention, once decided. Older recordings use
names chosen by each user, such as \file{2024-11-15-A/clicks}.}
